% math0 manuscript. Build with: latexmk -pdf paper.tex
\documentclass[11pt]{amsart}
\usepackage[margin=1.1in]{geometry}
\usepackage{amsmath,amssymb,amsthm,mathtools}
\usepackage{url}
\usepackage{hyperref}
\usepackage{xcolor}
\usepackage{listings}
\usepackage{newunicodechar}

\lstdefinelanguage{lean}{
  keywords={theorem,lemma,def,abbrev,example,fun,by,have,show,from,let,in,if,then,else,match,with,
    namespace,end,open,import,section,variable,instance,structure,where,calc,obtain,intro,exact,
    refine,rcases,rintro,simp,omega,decide,norm_num,linarith,nlinarith,positivity,induction,cases,
    private,noncomputable,unfold,rw,at,apply,ring,linear_combination},
  sensitive=true, comment=[l]{--}, morecomment=[s]{/-}{-/}, string=[b]"}
\lstset{language=lean, basicstyle=\ttfamily\footnotesize, columns=fullflexible, keepspaces=true, breaklines=true,
  keywordstyle=\color{blue!60!black}\bfseries, commentstyle=\color{gray}, frame=single,
  extendedchars=true, inputencoding=utf8,
  literate={ℕ}{{$\mathbb{N}$}}1 {ℤ}{{$\mathbb{Z}$}}1 {ℚ}{{$\mathbb{Q}$}}1 {∀}{{$\forall$}}1
    {↔}{{$\leftrightarrow$}}1 {¬}{{$\neg$}}1 {←}{{$\leftarrow$}}1 {→}{{$\to$}}1 {·}{{$\cdot$}}1
    {∑}{{$\sum$}}1 {∈}{{$\in$}}1 {∧}{{$\wedge$}}1 {×}{{$\times$}}1 {ˢ}{{$^{s}$}}1
    {⁻}{{$^{-}$}}1 {¹}{{$^{1}$}}1 {≤}{{$\le$}}1 {≠}{{$\ne$}}1 {₀}{{$_0$}}1
    {⟦}{{$[\![$}}1 {⟧}{{$]\!]$}}1}
% Unicode commonly found in Lean sources.
\newunicodechar{ℕ}{\ensuremath{\mathbb{N}}}
\newunicodechar{ℤ}{\ensuremath{\mathbb{Z}}}
\newunicodechar{ℚ}{\ensuremath{\mathbb{Q}}}
\newunicodechar{ℝ}{\ensuremath{\mathbb{R}}}
\newunicodechar{ℂ}{\ensuremath{\mathbb{C}}}
\newunicodechar{∀}{\ensuremath{\forall}}
\newunicodechar{∃}{\ensuremath{\exists}}
\newunicodechar{→}{\ensuremath{\to}}
\newunicodechar{←}{\ensuremath{\leftarrow}}
\newunicodechar{↔}{\ensuremath{\leftrightarrow}}
\newunicodechar{↦}{\ensuremath{\mapsto}}
\newunicodechar{¬}{\ensuremath{\neg}}
\newunicodechar{∧}{\ensuremath{\wedge}}
\newunicodechar{∨}{\ensuremath{\vee}}
\newunicodechar{≤}{\ensuremath{\le}}
\newunicodechar{≥}{\ensuremath{\ge}}
\newunicodechar{≠}{\ensuremath{\ne}}
\newunicodechar{∈}{\ensuremath{\in}}
\newunicodechar{∉}{\ensuremath{\notin}}
\newunicodechar{⊆}{\ensuremath{\subseteq}}
\newunicodechar{∩}{\ensuremath{\cap}}
\newunicodechar{∪}{\ensuremath{\cup}}
\newunicodechar{∑}{\ensuremath{\sum}}
\newunicodechar{∏}{\ensuremath{\prod}}
\newunicodechar{⟨}{\ensuremath{\langle}}
\newunicodechar{⟩}{\ensuremath{\rangle}}
\newunicodechar{·}{\ensuremath{\cdot}}
\newunicodechar{∘}{\ensuremath{\circ}}
\newunicodechar{⁻}{\ensuremath{^{-}}}
\newunicodechar{¹}{\ensuremath{^{1}}}
\newunicodechar{²}{\ensuremath{^{2}}}
\newunicodechar{₀}{\ensuremath{_{0}}}
\newunicodechar{₁}{\ensuremath{_{1}}}
\newunicodechar{₂}{\ensuremath{_{2}}}
\newunicodechar{α}{\ensuremath{\alpha}}
\newunicodechar{β}{\ensuremath{\beta}}
\newunicodechar{ε}{\ensuremath{\varepsilon}}
\newunicodechar{δ}{\ensuremath{\delta}}
\newunicodechar{π}{\ensuremath{\pi}}
\newunicodechar{σ}{\ensuremath{\sigma}}
\newunicodechar{φ}{\ensuremath{\varphi}}
\newunicodechar{∞}{\ensuremath{\infty}}
\newunicodechar{⌊}{\ensuremath{\lfloor}}
\newunicodechar{⌋}{\ensuremath{\rfloor}}
\newunicodechar{⌈}{\ensuremath{\lceil}}
\newunicodechar{⌉}{\ensuremath{\rceil}}
\newunicodechar{∣}{\ensuremath{\mid}}
\newunicodechar{≡}{\ensuremath{\equiv}}
\newunicodechar{×}{\ensuremath{\times}}
\newunicodechar{⧸}{\ensuremath{/}}
\newunicodechar{‖}{\ensuremath{\|}}
\newunicodechar{√}{\ensuremath{\surd}}
\newunicodechar{∅}{\ensuremath{\emptyset}}
\newunicodechar{⦃}{\ensuremath{\{\!|}}
\newunicodechar{⦄}{\ensuremath{|\!\}}}
\newunicodechar{𝓝}{\ensuremath{\mathcal{N}}}
\newunicodechar{≫}{\ensuremath{\gg}}
\newunicodechar{≪}{\ensuremath{\ll}}

\theoremstyle{plain}
\newtheorem{theorem}{Theorem}
\newtheorem{lemma}[theorem]{Lemma}
\newtheorem{proposition}[theorem]{Proposition}
\newtheorem{corollary}[theorem]{Corollary}
\newtheorem{conjecture}[theorem]{Conjecture}
\theoremstyle{definition}
\newtheorem{definition}[theorem]{Definition}
\newtheorem{remark}[theorem]{Remark}

\emergencystretch=3em
\raggedbottom
% Lean identifiers in running text (breakable at dots and underscores).
\expandafter\def\expandafter\UrlBreaks\expandafter{\UrlBreaks\do\_\do\.}
\DeclareUrlCommand\decl{\urlstyle{tt}}
\newcommand{\N}{\mathbb{N}}
\newcommand{\Q}{\mathbb{Q}}
\newcommand{\Z}{\mathbb{Z}}
\newcommand{\Qx}{\Q[\![X]\!]}

\title[Smallest sides of scalene integer triangles (OEIS A308107)]{The sum of the smallest sides of scalene integer triangles of perimeter $n$: a proof of Barker's generating function and recurrence for OEIS A308107}
\author{math0}
\address{Produced by Claude (an AI model developed by Anthropic) in the math0 project}
\date{October 24, 2026}

\begin{document}
\begin{abstract}
Let $a(n)$ be the sum of the smallest sides of all scalene triangles with integer sides and perimeter $n$ (OEIS A308107).
We prove the two formulas that Colin Barker conjectured for this sequence in 2019: the generating function $\sum_n a(n)x^n = x^9(2+2x+2x^2+x^3)/D(x)$, where $D(x)=(1-x)^4(1+x)^3(1+x^2)^2(1+x+x^2)^2$, and the resulting linear recurrence of order $15$, valid for $n>15$.
The proof separates the triangles whose middle side exceeds the smallest side by one, parametrizes these by two positive integers, and turns a chain of bijective shift relations into an identity of formal power series.
To our knowledge this is the first recorded proof (the entry still lists both formulas as conjectures), although the argument is elementary.
The result is verified in Lean~4 with Mathlib against the frozen statement that math0 formalized from the OEIS entry; that statement was reviewed for fidelity before any proof attempt, and its definition of $a$ is checked by the kernel against all $61$ listed terms.
\end{abstract}
\maketitle

\section{Introduction}

For $n\in\N$ let
\[
  T(n)=\{(k,i,j)\in\N^3 : k<i<j<k+i,\ k+i+j=n\},\qquad a(n)=\sum_{(k,i,j)\in T(n)} k .
\]
A scalene triangle with integer sides is determined by its three pairwise distinct side lengths, listed in increasing order as $k<i<j$; for such a triple the triangle inequalities reduce to $j<k+i$, and they force $k>j-i\ge1$.
So $T(n)$ is the set of scalene integer triangles of perimeter $n$, $k$ is the smallest side, and $a(n)$ is the sequence A308107 of the OEIS \cite{A308107}, submitted by Wesley Ivan Hurt on May~13, 2019, with offset $1$.
Its first nonzero terms are $a(9)=2$ (the triangle $(2,3,4)$), $a(11)=2$ (the triangle $(2,4,5)$) and $a(12)=3$ (the triangle $(3,4,5)$).
On the same day Colin Barker added two conjectured formulas to the entry (Conjecture~\ref{conj} below).
We prove both.

\begin{theorem}\label{thm:main}
As formal power series over $\Q$,
\begin{equation}\label{eq:gf}
  \sum_{n\ge0} a(n)\,x^n=\frac{x^9\,(2+2x+2x^2+x^3)}{(1-x)^4(1+x)^3(1+x^2)^2(1+x+x^2)^2},
\end{equation}
and for every integer $n>15$,
\begin{multline}\label{eq:rec}
  a(n) = -a(n-1) + 2a(n-3) + 4a(n-4) + 2a(n-5) - a(n-6) - 5a(n-7) - 5a(n-8) \\
  - a(n-9) + 2a(n-10) + 4a(n-11) + 2a(n-12) - a(n-14) - a(n-15).
\end{multline}
\end{theorem}

By the factorizations in the proof of Proposition~\ref{prop:AD}, the right-hand side of \eqref{eq:gf} equals
\[
  \frac{x^9\,(2-x^3-x^4)}{(1-x^2)(1-x^3)^2(1-x^4)^2}=\sum_{u,v,w\ge1}(v+w)\,x^{2u+4v+3w},
\]
which reflects the parametrization $(k,i,j)=(v+w,\,v+w+u,\,2v+w+u)$ of $T(n)$ by the positive integers $u,v,w$ with $2u+4v+3w=n$.
The Lean proof, which we follow in Section~\ref{sec:proof}, uses this parametrization in steps: it removes $u$ by the shift $(k,i,j)\mapsto(k,i+1,j+1)$, and then $v$ and $w$ by further shifts, each giving an identity $(1-X^d)F=G$ in $\Qx$.

\subsection*{Prior work and status}
The number of integer triangles of perimeter $n$ is Alcuin's sequence A005044 \cite{A005044}, with generating function $x^3/((1-x^2)(1-x^3)(1-x^4))$; Andrews, Paule and Riese \cite{APR01} extended this to $k$-gons with MacMahon's partition analysis, Hirschhorn \cite{Hir02} gave a short derivation of their result, and Patronis and Rizos \cite{PR23} compare two elementary methods for triangles.
These works count triangles and do not consider sums of side lengths.
Hurt's related entries A307686 (smallest sides, all triangles) \cite{A307686}, A308108 and A308109 (largest and middle sides, scalene) \cite{A308108,A308109} and A308158 (smallest sides, isosceles) \cite{A308158} also carry generating functions conjectured by Barker in May 2019.
These entries, and A308107 itself (revision \#14 of September~12, 2021, the current one), still mark Barker's formulas as conjectures and record no proof.
A comment of Roger Cuculière in A005044 relates the scalene triangles of perimeter $n+6$ to all triangles of perimeter $n$; this reduces A308107 to A307686 (Remark~\ref{rem:A307686}) but proves neither.
Fried's proofs of OEIS conjectures \cite{Fri24,Fri25,Fri26,Fri26b} do not cover these sequences.
The math0 literature referees found no prior proof or counterexample, so we classify Theorem~\ref{thm:main} as a new result; the argument is a routine computation of MacMahon type, and the formula may have been derived elsewhere without being recorded.

\section{Statement}\label{sec:statement}

\subsection*{The OEIS conjecture}
The FORMULA section of A308107 reads, verbatim:
\begin{lstlisting}[language={}]
Conjectures from _Colin Barker_, May 13 2019: (Start)
G.f.: x^9*(2 + 2*x + 2*x^2 + x^3) / ((1 - x)^4*(1 + x)^3*(1 + x^2)^2*(1 + x + x^2)^2).
a(n) = -a(n-1) + 2*a(n-3) + 4*a(n-4) + 2*a(n-5) - a(n-6) - 5*a(n-7) - 5*a(n-8) - a(n-9) + 2*a(n-10) + 4*a(n-11) + 2*a(n-12) - a(n-14) - a(n-15) for n>15.
(End)
\end{lstlisting}

\begin{conjecture}[Barker, 2019]\label{conj}
The sequence $a$ satisfies \eqref{eq:gf}, and \eqref{eq:rec} for all $n>15$.
\end{conjecture}

The entry's defining formula (also used in its PARI and Mathematica programs) is
\[
  a(n)=\sum_{k=1}^{\lfloor(n-1)/3\rfloor}\ \sum_{i=k+1}^{\lfloor(n-k-1)/2\rfloor}\operatorname{sign}\Bigl(\Bigl\lfloor\frac{i+k}{n-i-k+1}\Bigr\rfloor\Bigr)\,k .
\]
With $j=n-i-k$, the bound on $i$ says $i<j$, the sign factor is $1$ if $j<k+i$ and $0$ otherwise, and the bound on $k$ follows from $3k<k+i+j=n$; so this is the definition of $a$ above.

\subsection*{Lean statement}
The statement is the definition \decl{Math0.Statements.A308107.statement} in the math0 statement file \path{lean/Math0/Statements/OEIS/A308107.lean}, which was reviewed for fidelity and frozen (SHA-256 \texttt{e5f81c2e\ldots 516275}) before any proof attempt.
The definition of $a$ and the statement read, verbatim (docstrings omitted; namespace \decl{Math0.Statements.A308107}, with \texttt{open PowerSeries}):
\begin{lstlisting}
def a (n : ℕ) : ℕ :=
  ∑ t ∈ (Finset.range (n + 1) ×ˢ Finset.range (n + 1) ×ˢ Finset.range (n + 1)).filter
      (fun t : ℕ × ℕ × ℕ => t.1 < t.2.1 ∧ t.2.1 < t.2.2 ∧ t.2.2 < t.1 + t.2.1 ∧
        t.1 + t.2.1 + t.2.2 = n),
    t.1

def statement : Prop :=
  (PowerSeries.mk (fun n => (a n : ℚ)) =
    X ^ 9 * (2 + 2 * X + 2 * X ^ 2 + X ^ 3) *
      ((1 - X) ^ 4 * (1 + X) ^ 3 * (1 + X ^ 2) ^ 2 * (1 + X + X ^ 2) ^ 2)⁻¹) ∧
  (∀ n : ℕ, 15 < n →
    (a n : ℤ) = -(a (n - 1) : ℤ) + 2 * (a (n - 3) : ℤ) + 4 * (a (n - 4) : ℤ) +
      2 * (a (n - 5) : ℤ) - (a (n - 6) : ℤ) - 5 * (a (n - 7) : ℤ) - 5 * (a (n - 8) : ℤ) -
      (a (n - 9) : ℤ) + 2 * (a (n - 10) : ℤ) + 4 * (a (n - 11) : ℤ) + 2 * (a (n - 12) : ℤ) -
      (a (n - 14) : ℤ) - (a (n - 15) : ℤ))
\end{lstlisting}

\subsection*{Encoding}
A triple \texttt{t} stands for $(k,i,j)=(\texttt{t.1},\texttt{t.2.1},\texttt{t.2.2})$, so \texttt{a n} is $\sum_{T(n)}k$; the range $\{0,\dots,n\}^3$ loses nothing as the sides sum to $n$.
The entry has offset $1$; the extra value $a(0)=0$ (an empty sum) is harmless, as the right-hand side of \eqref{eq:gf} has zero constant term.
The first clause is \eqref{eq:gf} in Mathlib's ring $\Qx$, where \texttt{PowerSeries.mk f} is $\sum_n f(n)X^n$ and \texttt{⁻¹} is the inverse of power series over a field.
The denominator $D=(1-X)^4(1+X)^3(1+X^2)^2(1+X+X^2)^2$ has constant coefficient $1$, so $D^{-1}$ is its true inverse, and \decl{genFun_iff} (from Mathlib's \decl{PowerSeries.eq_mul_inv_iff_mul_eq}) shows that the first clause is equivalent to
\begin{equation}\label{eq:AD}
  A\cdot D = X^9(2+2X+2X^2+X^3),\qquad A=\textstyle\sum_{n\ge0}a(n)X^n .
\end{equation}
The second clause is \eqref{eq:rec} in $\Z$ for $n>15$; then $n-15\ge1$, so no truncated subtraction occurs.

\subsection*{Checks against the OEIS data}
The statement file defines a faster evaluator \decl{aFast}, which sums over the pairs $(k,i)$ with $j=n-(k+i)$, proves \decl{a_eq_aFast} (\texttt{a n = aFast n} for all $n$), and checks in \decl{a_data}, by \texttt{decide +kernel}, that $a(1),\dots,a(61)$ are the $61$ listed terms $0,\dots,0,2,0,2,3,5,3,9,\dots,740,735,825$.
Two further checks concern the transcription of the conjecture: \decl{recurrence_data} checks by \texttt{decide} that the listed terms satisfy \eqref{eq:rec} for $16\le n\le61$, and \decl{den_eq} checks by \texttt{ring} that
\begin{equation}\label{eq:den}
  D = 1 + X - 2X^3 - 4X^4 - 2X^5 + X^6 + 5X^7 + 5X^8 + X^9 - 2X^{10} - 4X^{11} - 2X^{12} + X^{14} + X^{15},
\end{equation}
the characteristic polynomial of \eqref{eq:rec}.

\section{Proof}\label{sec:proof}

We follow the declarations of the Lean file, whose names we give in brackets; all except \decl{result} are private.
We write $[P]$ for $1$ if $P$ holds and $0$ otherwise.
Besides $T(n)$ (\decl{T}; $a(n)=\sum_{T(n)}k$ is \decl{a_eq_T}) we use $U(n)=\{(v,w)\in\N^2 : v,w\ge1,\ 4v+3w+2=n\}$ (\decl{U}) and
\[
  b_v(n)=\sum_{U(n)}v,\quad b_w(n)=\sum_{U(n)}w,\quad c(n)=\#U(n),\quad e(n)=\#\{(v,w)\in U(n): v=1\}
\]
(\decl{bv}, \decl{bw}, \decl{c}, \decl{e}); in Lean, $T(n)$ and $U(n)$ are filtered from $\{0,\dots,n\}^3$ and $\{0,\dots,n\}^2$.
Each bijection below is applied with Mathlib's \decl{Finset.sum_nbij'}, which needs the two maps, membership of their images, the two inverse identities and equality of the summands; after unfolding membership, these goals are closed by \texttt{omega} (a few trivial ones by \texttt{simp} or \texttt{rfl}).

\subsection*{Step 1: the middle side}

\begin{lemma}[\decl{sum_T_filter}]\label{lem:tight}
For every $n$, $\displaystyle\sum_{(k,i,j)\in T(n),\ i=k+1}k=\sum_{(v,w)\in U(n)}(v+w)$.
\end{lemma}

\begin{proof}
Let $\varphi(v,w)=(v+w,\,v+w+1,\,2v+w+1)$ and $\psi(k,i,j)=(j-k-1,\,2k+1-j)$.
If $(v,w)\in U(n)$, then $(k,i,j)=\varphi(v,w)$ satisfies $i=k+1$, $k<i$, $i<j$ (as $v\ge1$), $j<k+i=2v+2w+1$ (as $w\ge1$) and $k+i+j=4v+3w+2=n$.
Conversely, if $(k,i,j)\in T(n)$ and $i=k+1$, then $v=j-k-1\ge1$ because $j>i$, $w=2k+1-j\ge1$ because $j<k+i$, and $4v+3w+2=j+2k+1=k+i+j=n$.
Direct computation gives $\varphi(\psi(k,i,j))=(k,i,j)$ when $i=k+1$, and $\psi(\varphi(v,w))=(v,w)$.
The summand $k$ at $\varphi(v,w)$ is $v+w$.
\end{proof}

\begin{lemma}[\decl{sum_T_filter_not}]\label{lem:loose}
For every $n$, $\displaystyle\sum_{(k,i,j)\in T(n+2),\ i\ne k+1}k=\sum_{(k,i,j)\in T(n)}k$.
\end{lemma}

\begin{proof}
The map $(k,i,j)\mapsto(k,i+1,j+1)$ sends $T(n)$ into $\{(k,i,j)\in T(n+2): i\ne k+1\}$, as $k<i$ gives $i+1\ne k+1$.
Its inverse is $(k,i,j)\mapsto(k,i-1,j-1)$: if $(k,i,j)\in T(n+2)$ and $i\ne k+1$, then $i\ge k+2$, so $k<i-1<j-1<k+(i-1)$ and the perimeter drops by $2$.
Both maps preserve the summand $k$.
\end{proof}

Splitting $T(n+2)$ according to whether $i=k+1$ and applying Lemmas~\ref{lem:tight} and~\ref{lem:loose} gives
\begin{equation}\label{eq:ashift}
  a(n+2)=a(n)+b_v(n+2)+b_w(n+2)\qquad(n\in\N)\quad(\text{\decl{a_shift}}).
\end{equation}

\subsection*{Step 2: shifts in \texorpdfstring{$U$}{U}}

\begin{lemma}[\decl{U_shift_v}, \decl{U_shift_w}]\label{lem:Ushift}
For every $n$ and every $g:\N^2\to\N$,
\begin{align*}
  \sum_{U(n+4)}g(v,w)&=\sum_{U(n)}g(v+1,w)+\sum_{U(n+4)}[v=1]\,g(v,w),\\
  \sum_{U(n+3)}g(v,w)&=\sum_{U(n)}g(v,w+1)+\sum_{U(n+3)}[w=1]\,g(v,w).
\end{align*}
\end{lemma}

\begin{proof}
The map $(v,w)\mapsto(v+1,w)$ is a bijection from $U(n)$ onto $\{(v,w)\in U(n+4): v\ne1\}$, with inverse $(v,w)\mapsto(v-1,w)$, since $4(v+1)+3w+2=(4v+3w+2)+4$.
Splitting $U(n+4)$ according to whether $v=1$ gives the first identity.
The second is the same argument with $(v,w)\mapsto(v,w+1)$, which raises $4v+3w+2$ by $3$.
\end{proof}

\begin{lemma}[\decl{sum_U_small}]\label{lem:small}
$U(n)=\emptyset$ for $n<9$; hence $b_v(n)=b_w(n)=c(n)=e(n)=0$ for $n<9$.
\end{lemma}

\begin{proof}
If $v,w\ge1$ then $4v+3w+2\ge9$.
\end{proof}

\begin{lemma}[\decl{bv_shift}, \decl{bw_shift}, \decl{c_shift}, \decl{e_shift}]\label{lem:rels}
For every $n\in\N$:
\begin{align*}
&(1)\ \ b_v(n+4)=b_v(n)+c(n+4), && (2)\ \ b_w(n+3)=b_w(n)+c(n+3),\\
&(3)\ \ c(n+4)=c(n)+e(n+4), && (4)\ \ e(n+3)=e(n)+[n+3=9].
\end{align*}
\end{lemma}

\begin{proof}
(3) is the first identity of Lemma~\ref{lem:Ushift} with $g=1$.
For (1), the first identity with $g(v,w)=v$ gives $b_v(n+4)=b_v(n)+c(n)+\sum_{U(n+4)}[v=1]\,v$; the last sum equals $e(n+4)$ because $v=1$ in each nonzero term, and $c(n)+e(n+4)=c(n+4)$ by (3).
For (2), the second identity with $g(v,w)=w$ gives $b_w(n+3)=b_w(n)+c(n)+\#\{(v,w)\in U(n+3):w=1\}$ in the same way, and the second identity with $g=1$ shows that $c(n)$ plus this count is $c(n+3)$.
For (4), the second identity with $g(v,w)=[v=1]$ gives $e(n+3)=e(n)+\sum_{U(n+3)}[w=1][v=1]$, and the last sum is $[(1,1)\in U(n+3)]=[n+3=9]$ because $4\cdot1+3\cdot1+2=9$.
\end{proof}

\subsection*{Step 3: power series}
We write $\mathrm{mk}\,f=\sum_{n\ge0}f(n)X^n\in\Qx$ for $f:\N\to\Q$, as in Mathlib.

\begin{lemma}[\decl{ps_shift}]\label{lem:ps}
Let $d\in\N$ and $f,g:\N\to\Q$ with $f(n+d)=f(n)+g(n+d)$ for all $n$ and $f(n)=g(n)$ for $n<d$.
Then $(1-X^d)\cdot\mathrm{mk}\,f=\mathrm{mk}\,g$.
\end{lemma}

\begin{proof}
The coefficient of $X^n$ on the left is $f(n)-f(n-d)$ if $n\ge d$ and $f(n)$ if $n<d$ (Mathlib's \decl{PowerSeries.coeff_X_pow_mul'}).
In the first case write $n=m+d$; then $f(m+d)-f(m)=g(m+d)$ by hypothesis. In the second case $f(n)=g(n)$.
\end{proof}

\begin{proposition}[\decl{gf_mul}]\label{prop:AD}
With $A=\mathrm{mk}\,a$ and $D$ as in \eqref{eq:AD}, $A\cdot D=X^9(2+2X+2X^2+X^3)$ in $\Qx$.
\end{proposition}

\begin{proof}
Let $B_v,B_w,C,E$ be the power series with coefficients $b_v,b_w,c,e$ (cast to $\Q$).
We apply Lemma~\ref{lem:ps} five times:
\begin{align*}
  (1-X^2)A&=B_v+B_w && \text{by \eqref{eq:ashift}, $d=2$},\\
  (1-X^4)B_v&=C && \text{by Lemma~\ref{lem:rels}(1), $d=4$},\\
  (1-X^3)B_w&=C && \text{by Lemma~\ref{lem:rels}(2), $d=3$},\\
  (1-X^4)C&=E && \text{by Lemma~\ref{lem:rels}(3), $d=4$},\\
  (1-X^3)E&=X^9 && \text{by Lemma~\ref{lem:rels}(4), $d=3$, with $g(n)=[n=9]$}.
\end{align*}
The initial conditions $f(n)=g(n)$ for $n<d$ hold because $b_v,b_w,c,e$ vanish for $n<9$ (Lemma~\ref{lem:small}), $a(0)=a(1)=0$ (\decl{a_zero_one}, checked by \texttt{decide}), and $[n=9]=0$ for $n<3$; the first relation also uses $\mathrm{mk}(b_v+b_w)=B_v+B_w$.
Put $P=1-X^2$, $Q=1-X^3$, $R=1-X^4$. Then $PQR\,A=QR\,B_v+QR\,B_w=(Q+R)\,C$, and multiplying by $QR$,
\[
  PQ^2R^2A=(Q+R)\,QR\,C=(Q+R)\,Q\,E=(Q+R)X^9=X^9(2-X^3-X^4).
\]
Now $PQ^2R^2=(1-X)D$ and $2-X^3-X^4=(1-X)(2+2X+2X^2+X^3)$, so $(1-X)(A\cdot D)=(1-X)X^9(2+2X+2X^2+X^3)$.
Since $\Qx$ is an integral domain and $1-X\ne0$ (its constant coefficient is $1$), we may cancel $1-X$.
In Lean, the cancellation is \texttt{mul\_left\_cancel}$_0$, and the remaining identity is a single \texttt{linear\_combination} of the six relations with coefficients $Q^2R^2$, $Q^2R^2$, $Q^2R$, $QR^2$, $Q(Q+R)$ and $Q+R$, checked by \texttt{ring}.
\end{proof}

\begin{proof}[Proof of Theorem~\ref{thm:main} (\decl{result})]
The generating-function clause follows from Proposition~\ref{prop:AD} and the equivalence \decl{genFun_iff} of Section~\ref{sec:statement}.
For the recurrence, write $D=\sum_{m=0}^{15}d_mX^m$ as in \eqref{eq:den}, so $d_0=1$, and let $n>15$.
The coefficient of $X^n$ in $X^mA$ is $a(n-m)$ for $m\le n$ (\decl{coeff_shift}), so the coefficient of $X^n$ in $A\cdot D$ is $\sum_{m=0}^{15}d_m\,a(n-m)$; in $X^9(2+2X+2X^2+X^3)=2X^9+2X^{10}+2X^{11}+X^{12}$ it is $0$, as $n\notin\{9,\dots,12\}$.
Hence $a(n)=-\sum_{m=1}^{15}d_m\,a(n-m)$ in $\Q$, which is \eqref{eq:rec} by \eqref{eq:den}; as the cast $\Z\to\Q$ is injective (\decl{Int.cast_injective}), \eqref{eq:rec} holds in $\Z$.
\end{proof}

\begin{remark}[A307686; not formalized]\label{rem:A307686}
Let $s(n)$ be the sum of the smallest sides of all integer triangles $k\le i\le j<k+i$ of perimeter $n$ (A307686) and $t(n)$ their number (A005044).
Cuculière's correspondence $(k,i,j)\mapsto(k-1,i-2,j-3)$ is a bijection from the scalene triangles of perimeter $n+6$ onto all triangles of perimeter $n$, and it lowers the smallest side by $1$; hence $a(n+6)=s(n)+t(n)$ for all $n\ge0$.
Since $a(n)=0$ for $n<6$ and $\sum t(n)x^n=x^3/((1-x^2)(1-x^3)(1-x^4))=x^3(1+x+x^2)(1-x^4)/D(x)$ (see \cite{APR01,Hir02}), Theorem~\ref{thm:main} gives
\[
  \sum_{n\ge0}s(n)x^n=\frac{x^3\bigl(2+2x+2x^2+x^3-(1+x+x^2)(1-x^4)\bigr)}{D(x)}=\frac{x^3\sum_{m=0}^{6}x^m}{D(x)},
\]
which is Barker's conjectured generating function in A307686. This is not part of the Lean verification; we checked $a(n+6)=s(n)+t(n)$ numerically for $n<194$.
\end{remark}

\section{Machine verification}\label{sec:lean}

\subsection*{Files and declarations}
The proof is in the file \path{lean/Math0/Frontier/OEIS/A308107.lean} of the math0 repository, \url{https://github.com/Foundation-0/math0}.
It imports the frozen statement file \path{lean/Math0/Statements/OEIS/A308107.lean} (Section~\ref{sec:statement}) and the audit commands in \path{lean/Math0/Audit.lean}.
The main declaration is \decl{Math0.Frontier.A308107.result}, stated as \texttt{theorem result : statement}; the private auxiliary declarations are those named in Section~\ref{sec:proof}, for instance
\begin{lstlisting}
private theorem gf_mul :
    mk (fun n => (a n : ℚ)) * den = X ^ 9 * (2 + 2 * X + 2 * X ^ 2 + X ^ 3)
\end{lstlisting}
The file uses no \texttt{sorry}, \texttt{native\_decide}, \texttt{set\_option} or macros, and its only kernel evaluation is \decl{a_zero_one} ($a(0)=a(1)=0$, by \texttt{decide}); the computational checks of the data, \decl{a_data} and \decl{recurrence_data}, are in the statement file, where \texttt{decide +kernel} makes the kernel evaluate \decl{aFast} on $n\le61$ (about $8\cdot10^4$ pairs in total) and compare with the listed terms.

\subsection*{Audit}
The proof file ends with \texttt{\#math0\_settles} applied to \decl{result} and \decl{statement}. This command, defined in \path{lean/Math0/Audit.lean} for statements that math0 formalizes itself, fails unless \decl{statement} is a \texttt{Prop}-valued definition without universe parameters, both declarations use no axioms besides \decl{propext}, \decl{Classical.choice} and \decl{Quot.sound} (which rejects \decl{sorryAx}, the axiom \decl{Lean.ofReduceBool} behind \texttt{native\_decide}, and custom axioms), and the type of \decl{result} is syntactically the statement or its negation, up to metadata, with no unfolding.
Running \texttt{lake env lean Math0/Frontier/OEIS/A308107.lean} in the directory \path{lean} prints\par\nobreak\noindent
\begin{minipage}{\linewidth}
\begin{lstlisting}[basicstyle=\ttfamily\scriptsize]
math0 audit: 'Math0.Statements.A308107.statement' depends on axioms [Classical.choice, Quot.sound, propext]
math0 audit: 'Math0.Frontier.A308107.result' depends on axioms [Classical.choice, Quot.sound, propext]
math0 statement check: 'Math0.Frontier.A308107.result' proves 'Math0.Statements.A308107.statement'
\end{lstlisting}
\end{minipage}
so \decl{result} proves the statement and depends only on the three permitted axioms.

\subsection*{Toolchain and reproduction}
We use Lean~4 \cite{Lean4} \texttt{v4.33.1} and Mathlib \cite{Mathlib} \texttt{v4.33.1}, pinned in \path{lean/}; the repository also pins formal-conjectures \cite{FC} at commit \texttt{c77fdd4}, which does not contain A308107 and is not used by this result.
To reproduce, run \path{scripts/verify.sh} from the repository root; it fetches the Mathlib cache, builds all math0 results with \texttt{lake build} (re-running every audit command) and checks that the pinned \texttt{formal\_conjectures} and \texttt{mathlib} checkouts are unmodified.
For this result alone, run \texttt{lake build Math0.Statements.OEIS.A308107} and \texttt{lake env lean Math0/Frontier/OEIS/A308107.lean} in \path{lean}; the second takes under a minute.

{\small
\begin{thebibliography}{99}

\bibitem{APR01}
G.~E. Andrews, P.~Paule, and A.~Riese, \emph{MacMahon's partition analysis IX: $k$-gon partitions}, Bull. Austral. Math. Soc. \textbf{64} (2001), no.~2, 321--329.

\bibitem{FC}
M.~Firsching, P.~Lezeau, S.~Mercuri, M.~Z. Horv\'ath, Y.~Dillies, C.~S\"onne, E.~Wieser, F.~Zhang, T.~Hubert, B.~Ag\"uera y Arcas, and P.~Kohli, \emph{Formal Conjectures: An Open and Evolving Benchmark for Verified Discovery in Mathematics}, arXiv:2605.13171, 2026. Repository: \url{https://github.com/google-deepmind/formal-conjectures}.

\bibitem{Fri24}
S.~Fried, \emph{Proofs of some conjectures from the OEIS}, arXiv:2410.07237, 2024.

\bibitem{Fri25}
S.~Fried, \emph{Proofs of several conjectures from the OEIS}, J. Integer Seq. \textbf{28} (2025), Article 25.4.3.

\bibitem{Fri26}
S.~Fried, \emph{Proofs of ten conjectures from the OEIS}, J. Integer Seq. \textbf{29} (2026), Article 26.1.8.

\bibitem{Fri26b}
S.~Fried, \emph{Further proofs of conjectures from the OEIS}, arXiv:2607.24832, 2026.

\bibitem{Hir02}
M.~D. Hirschhorn, \emph{$k$-gon partitions}, Bull. Austral. Math. Soc. \textbf{66} (2002), no.~1, 149--150.

\bibitem{Mathlib}
The mathlib Community, \emph{The Lean mathematical library}, in: Proceedings of the 9th ACM SIGPLAN International Conference on Certified Programs and Proofs (CPP 2020), ACM, 2020, 367--381.

\bibitem{Lean4}
L.~de Moura and S.~Ullrich, \emph{The Lean 4 theorem prover and programming language}, in: Automated Deduction -- CADE 28, Lecture Notes in Comput. Sci. \textbf{12699}, Springer, 2021, 625--635.

\bibitem{A005044}
The OEIS Foundation Inc., \url{https://oeis.org/A005044} (Alcuin's sequence), entry by R.~G. Wilson~v.

\bibitem{A307686}
The OEIS Foundation Inc., \url{https://oeis.org/A307686}, entry by W.~I. Hurt, May~11, 2019.

\bibitem{A308107}
The OEIS Foundation Inc., \url{https://oeis.org/A308107}, entry by W.~I. Hurt, May~13, 2019; revision \#14, Sep.~12, 2021.

\bibitem{A308108}
The OEIS Foundation Inc., \url{https://oeis.org/A308108}, entry by W.~I. Hurt, May~13, 2019.

\bibitem{A308109}
The OEIS Foundation Inc., \url{https://oeis.org/A308109}, entry by W.~I. Hurt, May~13, 2019.

\bibitem{A308158}
The OEIS Foundation Inc., \url{https://oeis.org/A308158}, entry by W.~I. Hurt, May~14, 2019.

\bibitem{PR23}
T.~Patronis and I.~Rizos, \emph{Triangles with integer sides and given perimeter}, Mathematical Review \textbf{97--98} (2022), 78--89 (in Greek); arXiv:2308.01912.

\end{thebibliography}
}

\end{document}
